\documentclass[12pt]{article}
\usepackage[a4paper,margin=1in]{geometry}
\usepackage[utf8]{inputenc}
\usepackage[T1]{fontenc}
\usepackage{amsmath,amssymb}
\usepackage{textcomp}
\usepackage{parskip}
\begin{document}
\begin{center}
{\Large \textbf{2008 STEP III}}
\end{center}

\subsection*{Page 1}
\begin{flushleft}
\bigskip\textbf{Section A: Pure Mathematics}\par
\bigskip\textbf{Question 1.} Find all values of a, b, x and y that satisfy the simultaneous equations\par
\[a + b = 1\]
\[ax + by = 1\]
\[ax2 + by2 = 1\]
\[ax3 + by3 = 1 .\]
[ Hint: you may wish to start by multiplying the second equation by x + y. ]\par
n\par
\bigskip\textbf{Question 2.} Let S (n) rk , where k is a positive integer, so that\par
k\par
\equiv\par
\[r=0\]
S (n) 1 n(n + 1) and S (n) 1 n(n + 1)(2n + 1) .\par
1 \equiv 2 2 \equiv 6\par
n\par
\textbf{(i)} By considering (r + 1)k rk , show that\par
-\par
\[r=0\]
k k k\par
\[kS (n) = (n + 1)k (n + 1) S (n) S (n) S (n) . ( )\]
k 1 k 2 k 3 1\par
- - - 2 - - 3 - - \cdots - k 1 ∗\par
-\par
Obtain simplified expressions for S (n) and S (n).\par
3 4\par
\textbf{(ii)} Explain, using ( ), why S (n) is a polynomial of degree k + 1 in n. Show that in this\par
k\par
∗\par
polynomial the constant term is zero and the sum of the coefficients is 1.\par
\bigskip\textbf{Question 3.} The point P (a cos θ , b sin θ), where a > b > 0, lies on the ellipse\par
\[x^2 y^2\]
\[+ = 1 .\]
\[a^2 b^2\]
The point S( ea , 0), where b2 = a2(1 e2) , is a focus of the ellipse. The point N is the foot\par
- -\par
of the perpendicular from the origin, O, to the tangent to the ellipse at P . The lines SP and\par
ON intersect at T . Show that the y-coordinate of T is\par
b sin θ\par
.\par
1 + e cos θ\par
Show that T lies on the circle with centre S and radius a.\par
\par
  \par
     \par
\end{flushleft}
\subsection*{Page 2}
\begin{flushleft}
\bigskip\textbf{Question 4.} (i) Show, with the aid of a sketch, that y > tanh(y/2) for y > 0 and deduce that\par
x 1\par
\[arcosh x > - for x > 1. ( )\]
\[\sqrt\{x^2\} 1 ∗\]
-\par
x 1\par
\textbf{(ii)} By integrating ( ), show that arcosh x > 2 - for x > 1.\par
\[∗ \sqrt\{x^2\} 1\]
-\par
\[\sqrt\{x^2\} 1\]
\textbf{(iii)} Show that arcosh x > 3 - for x > 1.\par
x + 2\par
[Note: arcosh x is another notation for cosh 1 x. ]\par
-\par
\bigskip\textbf{Question 5.} The functions T (x), for n = 0, 1, 2, . . . , satisfy the recurrence relation\par
n\par
\[T n+1 (x) 2xT n (x) + T n 1 (x) = 0 (n  1). ( )\]
- - ∗\par
Show by induction that\par
\[(T (x))2 T (x)T (x) = f(x) ,\]
n n 1 n+1\par
- -\par
\[where f(x) = (T (x))2 T (x)T (x) .\]
1 0 2\par
-\par
In the case f(x) 0, determine (with proof) an expression for T (x) in terms of T (x) (assumed\par
n 0\par
\equiv\par
\[to be non-zero) and r(x), where r(x) = T (x)/T (x). Find the two possible expressions for\]
1 0\par
r(x) in terms of x.\par
dy\par
\bigskip\textbf{Question 6.} In this question, p denotes .\par
dx\par
\textbf{(i)} Given that\par
\[y = p^2 + 2xp ,\]
show by differentiating with respect to x that\par
dx 2x\par
\[= 2 .\]
dp - - p\par
\[Hence show that x = 2 p + Ap 2 , where A is an arbitrary constant.\]
- 3 -\par
\[Find y in terms of x if p = 3 when x = 2.\]
-\par
\textbf{(ii)} Given instead that\par
\[y = 2xp + p ln p ,\]
\[and that p = 1 when x = 1 , show that x = 1 ln p 1 and find y in terms of x.\]
- 4 - 2 - 4\par
[Turn over\par
\end{flushleft}
\subsection*{Page 3}
\begin{flushleft}
\bigskip\textbf{Question 7.} The points A, B and C in the Argand diagram are the vertices of an equilateral triangle\par
described anticlockwise. Show that the complex numbers a, b and c representing A, B and C\par
satisfy\par
\[2c = (a + b) + i\sqrt\{3\}(b a).\]
-\par
Find a similar relation in the case that A, B and C are the vertices of an equilateral triangle\par
described clockwise.\par
\textbf{(i)} The quadrilateral DEF G lies in the Argand diagram. Show that points P , Q, R and S\par
can be chosen so that P DE, QEF , RF G and SGD are equilateral triangles and P QRS\par
is a parallelogram.\par
\textbf{(ii)} The triangle LM N lies in the Argand diagram. Show that the centroids U , V and W\par
of the equilateral triangles drawn externally on the sides of LM N are the vertices of an\par
equilateral triangle.\par
[Note: The centroid of a triangle with vertices represented by the complex numbers\par
x, y and z is the point represented by 1 (x + y + z) .]\par
\bigskip\textbf{Question 8.} (i) The coefficients in the series\par
\[S = 1 x + 1 x^2 + 1 x^3 + + a xr +\]
3 6 12 \cdots r \cdots\par
\[satisfy a recurrence relation of the form a + pa = 0. Write down the value of p.\]
r+1 r\par
By considering (1 + px)S, find an expression for the sum to infinity of S (assuming that\par
it exists). Find also an expression for the sum of the first n + 1 terms of S.\par
\textbf{(ii)} The coefficients in the series\par
\[T = 2 + 8x + 18x2 + 37x3 + + a xr +\]
r\par
\cdots \cdots\par
\[satisfy a recurrence relation of the form a + pa + qa = 0. Find an expression for\]
r+2 r+1 r\par
the sum to infinity of T (assuming that it exists). By expressing T in partial fractions,\par
or otherwise, find an expression for the sum of the first n + 1 terms of T .\par
\end{flushleft}
\subsection*{Page 4}
\begin{flushleft}
\bigskip\textbf{Section B: Mechanics}\par
\bigskip\textbf{Question 9.} A particle of mass m is initially at rest on a rough horizontal surface. The particle experiences\par
a force mg sin \pit, where t is time, acting in a fixed horizontal direction. The coefficient of\par
friction between the particle and the surface is µ. Given that the particle starts to move first\par
\[at t = T , state the relation between T and µ.\]
0 0\par
\textbf{(i)} For µ = µ , the particle comes to rest for the first time at t = 1. Sketch the acceleration-\par
time graph for 0  t  1. Show that\par
\[1 + 1 µ2 2 µ \pi + µ arcsin µ = 0 .\]
- 0 - 0 0 0\par
\textbf{(ii)} For µ = µ 0 sketch the acceleration-time graph for 0  t  3. Describe the motion of the\par
\[particle in this case and in the case µ = 0.\]
[Note: arcsin x is another notation for sin 1 x. ]\par
-\par
\bigskip\textbf{Question 10.} A long string consists of n short light strings joined together, each of natural length  and\par
modulus of elasticity λ. It hangs vertically at rest, suspended from one end. Each of the short\par
strings has a particle of mass m attached to its lower end. The short strings are numbered 1\par
to n, the nth short string being at the top. By considering the tension in the rth short string,\par
determine the length of the long string. Find also the elastic energy stored in the long string.\par
A uniform heavy rope of mass M and natural length L has modulus of elasticity λ. The rope\par
hangs vertically at rest, suspended from one end. Show that the length, L, of the rope is given\par
by\par
M g\par
\[L = L 1 + ,\]
2λ\par
and find an expression in terms of L, L and λ for the elastic energy stored in the rope.\par
[Turn over\par
� \par
 \par
\end{flushleft}
\subsection*{Page 5}
\begin{flushleft}
\bigskip\textbf{Question 11.} A circular wheel of radius r has moment of inertia I about its axle, which is fixed in a horizontal\par
position. A light string is wrapped around the circumference of the wheel and a particle of\par
mass m hangs from the free end. The system is released from rest and the particle descends.\par
The string does not slip on the wheel.\par
As the particle descends, the wheel turns through n revolutions, and the string then detaches\par
from the wheel. At this moment, the angular speed of the wheel is ω . The wheel then turns\par
through a further n revolutions, in time T , before coming to rest. The couple on the wheel\par
due to resistance is constant.\par
Show that\par
\[ω T = 2\pin\]
0 2\par
and\par
\[mgrn T 2 4\pimr2n2\]
\[I = 1 - 2 .\]
\[4\pin (n + n )\]
2 1 2\par
\end{flushleft}
\subsection*{Page 6}
\begin{flushleft}
\bigskip\textbf{Section C: Probability and Statistics}\par
\bigskip\textbf{Question 12.} Let X be a random variable with a Laplace distribution, so that its probability density function\par
is given by\par
\[f(x) = e x , < x < . ( )\]
-| |\par
2 -\infty \infty ∗\par
\[Sketch f(x). Show that its moment generating function M (θ) is given by M (θ) = (1 θ2) 1\]
X X -\par
-\par
and hence find the variance of X.\par
A frog is jumping up and down, attempting to land on the same spot each time. In fact, in\par
each of n successive jumps he always lands on a fixed straight line but when he lands from the\par
ith jump (i = 1 , 2 , . . . , n) his displacement from the point from which he jumped is X cm,\par
i\par
where X has the distribution ( ). His displacement from his starting point after n jumps is\par
i\par
n ∗\par
\[Y cm (so that Y = X ). Each jump is independent of the others.\]
i\par
\[i=1\]
Obtain the moment generating function for Y /\sqrt2n and, by considering its logarithm, show\par
that this moment generating function tends to exp( 1 θ2) as n .\par
2 \to \infty\par
Given that exp( 1 θ2) is the moment generating function of the standard Normal random vari-\par
able, estimate the least number of jumps such that there is a 5\% chance that the frog lands\par
25 cm or more from his starting point.\par
\bigskip\textbf{Question 13.} A box contains n pieces of string, each of which has two ends. I select two string ends at\par
random and tie them together. This creates either a ring (if the two ends are from the same\par
string) or a longer piece of string. I repeat the process of tying together string ends chosen at\par
random until there are none left.\par
Find the expected number of rings created at the first step and hence obtain an expression\par
for the expected number of rings created by the end of the process. Find also an expression\par
for the variance of the number of rings created.\par
Given that ln 20 3 and that 1 + 1 + + 1 ln n for large n, determine approximately the\par
\approx 2 \cdots n \approx\par
\[expected number of rings created in the case n = 40 000.\]
\par
\end{flushleft}
\end{document}
